\section{Clustering}
This class is a type fo \textit{unsupervised learning} in which we assume that the 
data has some kind of unkwon structure, and we wan to figure it out. To do that we
try to accomplish two things: 
\begin{itemize}
    \item minimize distance between datapoints of the same cluster
    \item maximize distance between \diff clusters
\end{itemize}
\begin{figure}[h]
    \centering
    \includegraphics[width=0.3\textwidth]{clustering/cluster_example.png}    
\end{figure}
The clustering output is really subjective and depend what we're trying to achieve
and our conception of \textbf{similarity}. We have to define two elements:
\begin{description}
    \item[Distance metric]: tells what makes two object similar/dissimilar
    \item[Clustering algorithm]: procedure to compute groups/clusters
\end{description}
\vspace{0.5cm}

%-----------------------------------------------------------------------------------------
\subsubsection*{Distance metrics}
The definition of Distance Metrics depends mostly on the type of the data features.
Some common ones are:
\begin{itemize}
    \item \textbf{Euclidean distance}: $d_E(x,q)=\sqrt{\sum_{i=1}^n(x_1-q_1)^2}$

    \item \textbf{Squared Euclidean Distance}: $d_{E^2}(x,q)=(d_E(x,q))^2 
    \rightarrow$ with respect to $d_E$ it gives more emphasys to outliners 
    \textcolor{gray}{(distance is squared)}

    \item \textbf{Standardize Euclidean Distance}: $d_{SE}(x,q)=\sqrt{\frac{1}
    {s_i^2}(x_i-q_i)} \rightarrow$ weights each dimension by something inverse
    to its variability, favouring those dimensions that doens't change too much

    \item \textbf{Chebychev Distance}: $d_ {\max}(x,q)=\max_i|x_i-q_i|$ 
    \textcolor{gray}{(just the max, not really useful)}

    \item \textbf{Manhattan Distance}: $d_M(x,q)=\sum_{i=1}^n|x_i-q_i| \rightarrow$
    gives the \# step to reach a destination in a grid by moving up and down, lot of
    usage in many fields

    \item \textbf{Correlation-based Distance}: $d_R(x,q)=1-\cfrac{s_{x,q}}{\sqrt{s_x}
    \sqrt{s_q}}=1-\cfrac{\sum_i(x_i-\bar{x})(q_i-\bar{q})}{\sqrt{\sum_i(x_i-\bar{x})
    ^2\sum_i(q_i-\bar{q})^2}} \rightarrow$ this measures allow to get an idea about
    the distance between feature trends distance 
\end{itemize}
\hfill \\
Highly suggested to use $d_R$ in the project\\
\hfill \\
To compare the datapoint features, we must be sure that everyone is using the same
measurement unit.

\newpage
%-----------------------------------------------------------------------------------------
\subsubsection*{Clustering algorithms}
There are several way to categorize these methods:
\begin{description}
    \item[Hard clustering]: each datapoint can be only in one cluster
    \item[Soft clustering]: each datapont might be in more clusters at the same time
\end{description}
\begin{figure}[h]
    \centering
    \includegraphics[width=0.6\textwidth]{clustering/soft_vs_hard_clustering.png}
\end{figure}

\begin{description}
    \item[Flat clustering]: all clusters are at the same level
    \item[Hierachical clustering]: the cluster are under some hierarchy (dendogram)
\end{description}

% this is just a summary, place it at the begging of the relative section
\begin{comment}
Partition-based clustering: A partitioning method creates k partitions from given set of n data objects.
Initially, each data objects are assigned to some of the partitions. An iterative relocation technique is used to
improve the partitioning by moving objects from one group to another. E.g. k-means, k-median;

Hierarchical clustering: hierarchical clustering is a method of cluster analysis that seeks to build a hierarchy of
clusters. E.g. Agglomerative, Divisive;

• Density-based clustering: identify distinctive clusters in the data, based on the idea that a cluster/group in a
data space is a contiguous region of high point density, separated from other clusters by sparse regions. e.g.
DBSCAN;
is more similar ~

• Probabilistic clustering: models the data as generated by a mixture of probability distributions. Each
cluster is a component (e.g., a Gaussian) and each observation has a membership probability (soft
assignment)
not hard assignment, more flexibility 
\end{comment}

\vspace{0.5cm}
%-----------------------------------------------------------------------------------------
\subsection{K-means algorithm}
In this algorithm we priorly fix the number of clusters $k$, starting from an initial 
assignment and itereatively moving datapoints between the clusters.

% maybe make it like code snippet
\begin{enumerate}
    \item define $k \rightarrow C_1\dots C_k$ clusters
    \item randomly assign each point to some cluster $C_i$ and compute the centroid for 
    that set of points
    \item re-assign each point to the closest centroid cluster, ending up possibly with 
    new point in each cluster
    \item recompute the centroids, and iterate from \textit{step 2}
\end{enumerate} 
We iterate unitl the algorithm converges, e.g. the points won't change cluster.

\begin{figure}[h]
    \centering
    % --- First row ---
    \begin{subfigure}[t]{0.3\textwidth}
        \centering
        \includegraphics[width=\linewidth]{clustering/kmeans/kmean_step0.png}
    \end{subfigure}
    \hspace{0.3cm}
    \begin{subfigure}[t]{0.3\textwidth}
        \centering
        \includegraphics[width=\linewidth]{clustering/kmeans/kmean_step1.png}
    \end{subfigure}
    \hfill \\
    % --- Second row ---
    \begin{subfigure}[t]{0.3\textwidth}
        \centering
        \includegraphics[width=\linewidth]{clustering/kmeans/kmean_step2.png}
    \end{subfigure}
    \hspace{0.3cm}
    \begin{subfigure}[t]{0.3\textwidth}
        \centering
        \includegraphics[width=\linewidth]{clustering/kmeans/kmean_step3.png}
    \end{subfigure}
\end{figure}

\newpage
\textbf{Limitations} of this algorithm:
\begin{itemize}
    \item the algorithm output strongly depends on the initial random assignment; if 
    there is some structure probably the algorithm will converge at the same result, but 
    if this structure is not so evident the outcome might be very sensitive.\newline
    \begin{figure}[h]
        \centering
        \includegraphics[width=0.2\textwidth]{clustering/kmeans/init_assignment_error.png}
    \end{figure}
    To reduce this influence we could try out multiple seeds in parallel and choose at the 
    end the best result.

    \item we have to choose $k$ in advance, maybe without clear idea of we're doing
    
    \item the distance used affects a lot the centroid: if we use \textit{Euclidean 
    Distance} our computation is going to be very sensitive to outliers
    \begin{figure}[h]
        \centering
        \includegraphics[width=0.2\textwidth]{clustering/kmeans/k_medoids.png}
    \end{figure}
    We can resolve this by choosing another method to compute distances: instead of using 
    $d_E$ we use the \textbf{mode}, which are more robuts to utliers $\rightarrow$ the 
    \textit{medoid} is also more reasonable to use since we'll end up with something that 
    is a real point of the cluster \textcolor{gray}{(like the most representative point)}.
    
    \item with strange-shape clusters, K-means makes a very bad work (how can we know that 
    in advance? You can't) $\rightarrow$ k-means fits well spherical/hellipsoidal 
    datashapes
\end{itemize}

%-----------------------------------------------------------------------------------------
\subsection{Hierarchical clustering}
%It can be used for many many things, not only for clustering
This method builds a hierarchy among the clusters using a \textbf{Distance Matrix D}, which 
contains the distance between each pair of samples. To stop the procedure we have to define 
a \textbf{cutting criterion} to stop the construction of the \textbf{dendogram} := 
representation of your cluster.
\begin{figure}[h]
    \centering
    \includegraphics[width=0.5\textwidth]{clustering/hierachical/dendogram_example.png}
\end{figure}
The algorith can both start by considering each datapoint in a \diff cluster and then
merging them, or start from a one big cluster and progressively split it.

\begin{enumerate}
    \item start from considering each datapoints as a cluster per se
    \item merge the two closer clusters and take a "center" (linkage methods)
    \item recompute the set of points and iterate starting from 1
\end{enumerate}
By merging the cluster $D$ become smaller at each iteration.

By using this algorithm we can also highligh noisy datapoints and anomalies: if a point is 
added very late in the cluster, maybe it's just an anomaly $\rightarrow$ this method can
also be used for anomaly detection.

\paragraph{Linkage methods}\hfill\\
How do we computer the "center" between build-up clusters? 
\begin{description}
    \item[Single link]: cluster-distance is based on the two closest points from each 
    cluster
    \begin{figure}[h]
        \centering
        \includegraphics[width=0.2\textwidth]{clustering/hierachical/single_linkage.png}        
    \end{figure}

    \item[Average link]: cluster-distance is based on the average position of the cluster
    datapoints
    \begin{figure}[h]
        \centering
        \includegraphics[width=0.2\textwidth]{clustering/hierachical/avg_linkage.png}        
    \end{figure}

    \item[Complete link]: cluster-distance is based on the most distant points from each
    cluster
    \begin{figure}[h]
        \centering
        \includegraphics[width=0.2\textwidth]{clustering/hierachical/complete_linkage.png}        
    \end{figure}

\end{description}
There is no rule of thumb, start with the average and adjust things.
% what about the mode for this one?

%-----------------------------------------------------------------------------------------
\subsection{Evaluating cluster quality}
How do we eavulate if our clusters make sense if we don't have an output reference?

\subsubsection{Within Cluster Sum of Square (WCSS)}
WCSS measures the variability of the observations within each cluster.
\begin{equation}
    WCSS = \sum_i\sum_{p\in C_i} dist(p,C_i)^2
\end{equation}

The smaller our clusters are, the lower(better) the value of WSS. However by following
this metric, we get that the best result is we each datapoint form a cluster per-se. So we 
don't want to minimize this, otherwise is useless; we need just to find a point in which it
doesn't decreases a too much.\\
\hfill \\
With this method we're also looking at how similar points within the same cluster are, 
without considering the other goal of clustering: distance between \diff clusters.

\paragraph{Solhouette Score}
This metrics takes into account both the distance between points in the same cluster and
the distance from other cluster's points.

\begin{enumerate}
    \item $\forall$ point $i$ get all the distances from the points within the same
    cluster $a(i)$
    \item $\forall$ point $i$ get the distance from the points in the nearer cluster $b(i)$
    \item $\forall$ point $i$ get the average of both $\bar{a}(i),\bar{b}(i)$
    \item get the \textbf{Silhouette score} $\forall\ i$: $sil(i)=\cfrac{\bar{b}(i)-\bar{
    a}(i)}{\max\{\bar{a}(i),\bar{b}(i)\}}$
    \item get the \textbf{Silhouette score for cluster $k$}: $sil(C_k)=\frac{1}{|C_k|}\sum_{ 
    i\in C_k}sil(i)$
    \item get the \textbf{Silhouette score for overhall clustering}: $sil(C)=\cfrac{1}{K}
    \sum_{k=1}^K sil(C_k)$
\end{enumerate}